\section{Discussion}
\label{sec:discussion}

The experiments in Section~\ref{sec:results} deliver a result that
we did not expect at the outset. The closed-form rotation-plus-
scaling Jacobian of Equation~\eqref{eq:our-estimator}, despite being
exactly the backward signal that the forward quantizer performs, is
\emph{harmful} as a drop-in replacement for the STE at the
configuration we tested: it induces a sharp codebook collapse after
roughly $1{,}000$ optimisation steps and never recovers. The
decomposed ablation, Table~\ref{tab:ablation_modes}, localises the
damage to the component that interacts with the commitment loss. We
therefore make two separate claims that did not exist when we began.
First, the identity $\qvec = s\cdot\rmat\ze$ is a correct and
exploitable geometric fact; the Householder operator is an efficient
$\mathcal{O}(d)$ implementation of it. Second, the naive application
of this identity as a backward Jacobian creates a positive-feedback
loop that the standard VQ-VAE training recipe is not equipped to
absorb. A useful rotation-aware estimator therefore has to combine
the closed-form geometry with an explicit mechanism for preventing
the feedback loop --- for example, a codebook-entropy regulariser, a
temperature that anneals the rotation toward identity as training
progresses, or the use of only the rotation \emph{without} the
scalar rescaling. We leave a full characterisation of which
combination is stable to follow-up work and report the ablation
results in Section~\ref{sec:results}.

We now discuss limitations and one clarification.

\paragraph{Limitation: the forward operation is still hard.}
Our method changes only the backward pass; the forward quantizer
remains a bit-exact nearest-neighbour lookup. This is a deliberate
design choice because any downstream consumer that freezes the
tokenizer --- autoregressive transformers, discrete diffusion, masked
image modelling --- requires integer code indices.  It also means
our approach is not directly comparable to methods such as FSQ that
eliminate the codebook altogether, since those target a different
trade-off between representation simplicity and codebook
controllability.

\paragraph{Limitation: scale.}
The experiments reported here are on CIFAR-10 only. The $32 \times
32$ resolution permits a $12{,}000$-step comparison across four
baselines within a few GPU-hours and surfaces the codebook-collapse
regime reliably, but it leaves open whether the relative ordering of
methods persists at $128{\times}128$ or $256{\times}256$ resolution
with larger codebooks. Our software stack supports FFHQ and ImageNet
configurations; scaled experiments are the primary item on our
immediate follow-up list.

\paragraph{Limitation: downstream metrics.}
We evaluate tokenizers by their \emph{reconstruction} quality and
codebook utilisation, not yet by the generative-modelling quality of
a downstream autoregressive or diffusion model that uses those
codes. The correlation between tokenizer PSNR / perplexity and
final generation FID is well-documented~\citep{yu2022vitvqgan,
yu2023magvit2,yu2024titok}, but is not perfect. A full downstream
evaluation (training a MaskGIT-style transformer on top of each
tokenizer, measuring class-conditional FID) is planned as the next
experiment in our matrix (E5 in our experimental plan).

\paragraph{Clarification: relation to the rotation trick.}
\citet{fifty2024rotation} independently arrived at the idea that the
STE discards the rotation between $\ze$ and $\qvec$. Our technical
contribution relative to that work is the specific closed-form
Householder formulation of Section~\ref{sec:method}, the isolation
of the rotation and the rescaling components as independently
ablatable design choices, and the stop-gradient decomposition that
yields bit-exact forward output even in the degenerate
$\hat{\ze} \approx \hat{\qvec}$ regime. These are concrete
implementation choices, not grand theoretical differences; we
welcome direct comparison in future work and anticipate that a
combined estimator drawing on both lines could further improve the
quality-throughput Pareto.

\paragraph{Broader impact.}
Because our method is a drop-in replacement for the standard VQ
quantizer, its main effect is to make existing VQ-based
architectures slightly better. It does not enable new capabilities
that were previously infeasible, does not require new data
collection, and does not affect deployment-time compute.  The
ethical considerations of generative modelling more broadly ---
misuse for synthetic media, representational fairness of
tokenizers trained on web-scraped corpora --- apply to the larger
systems that would consume our tokenizer, not to the tokenizer
itself. We encourage downstream users to conduct domain-appropriate
audits before deployment.
